\BourbakiExercisesSection

\begin{enumerate}
\def\labelenumi{\arabic{enumi}.}
\tightlist
\item
  Let \(\alpha \mapsto f_{\alpha}\) be an action of \(\Omega\) on E. Let F be the image of \(\Omega\) in \(\mathbf{E}^{\mathbb{E}}\) and G the stable subset of \(\mathbf{E}^{\mathbb{E}}\) under the law \((f,g) \mapsto f \circ g\) generated by F and the identity mapping of E.
\end{enumerate}

\begin{enumerate}
\def\labelenumi{(\alph{enumi})}
\item
  Show that every subset of E which is stable under the action of \(\Omega\) is also stable under the restriction to G of the canonical action of \(E^{E}\) on E (no. 1, Example 3).
\item
  Let X be a subset of E; show that the stable subset of E generated by X is the set of \(f(x)\) where f runs through G and x runs through X.
\end{enumerate}

\begin{enumerate}
\def\labelenumi{\arabic{enumi}.}
\setcounter{enumi}{1}
\item
  Let \(\bot\) be an internal law on E which is doubly distributive with respect to the associative internal law \(\top\) ; show that, if \(x \perp x'\) and \(y \perp y'\) are cancellable under the law \(\top\) , \(x \perp y'\) is permutable with \(y \perp x'\) under the law \(\top\) (calculate the composition \((x \top y) \perp (x' \top y')\) in two different ways). In particular, if the law \(\bot\) has an identity element, two cancellable elements under the law \(\top\) are permutable under this law; if all the elements of E are cancellable under the law \(\top\) , this law is commutative.
\item
  Let \(\top\) and \(\bot\) be two internal laws on a set \(E\) such that \(\bot\) is right distributive with respect to \(\top\) .
\end{enumerate}

\begin{enumerate}
\def\labelenumi{(\alph{enumi})}
\item
  If the law \(\top\) has an identity element \(e\) , \(x \perp e\) is idempotent under the law \(\top\) for all \(x \in \mathbb{E}\) ; if further there exists \(y\) such that \(x \perp y\) is cancellable under the law \(\top\) , then \(x \perp e = e\) .
\item
  If the law \(\perp\) has an identity element u and there exists \(z \in E\) which is cancellable for both laws \(\perp\) and \(\top\) , u is cancellable under the law \(\top\) .
\end{enumerate}

¶ 4. Let \(\top\) and \(\bot\) be two internal laws on E, each with an identity element; if the left action of E on itself derived from each of the laws is distributive with respect to the other, every element of E is idempotent under both these laws (if \(e\) is the identity element for \(\top\) , \(u\) the identity element for \(\bot\) , prove first that \(e \perp e = e\) , noting that \(e = e \perp (u \top e)\) ).

\begin{enumerate}
\def\labelenumi{\arabic{enumi}.}
\setcounter{enumi}{4}
\tightlist
\item
  On a set E suppose three internal laws of composition are given: an addition (not necessarily associative nor commutative), a multiplication (not necessarily associative) and a law denoted by \(\top\) . Suppose that the multiplication has an identity element \(e\) , that the law \(\top\) is right distributive with respect to multiplication and the law \(\top\) is left distributive with respect to addition. Show that, if there exist \(x, y, z\) such that \(x \top z, y \top z\) and \((x + y) \top z\) are cancellable under multiplication, then \(e + e = e\) (use Exercise 3(a)).
\end{enumerate}

¶ 6. An internal law of composition \(\top\) on E is said to determine on E a quasi-group structure if, for all \(x \in E\) , the left and right translations \(\gamma_x\) and \(\delta_x\) are bijective mappings of E onto itself. A quasi-group is said to be distributive if the law \(\top\) is doubly distributive with respect to itself.

\begin{enumerate}
\def\labelenumi{(\alph{enumi})}
\item
  Determine all the distributive quasi-group structures on a set of \(n\) elements for \(2 \leqslant n \leqslant 6\) .
\item
  *Show that the set \(\mathbf{Q}\) of rational numbers, with the law of composition \((x,y)\mapsto\frac{1}{2}(x+y)\) , is a distributive quasi-group.*
\item
  In a distributive quasi-group E, every element is idempotent. Deduce that, if E has more than one element, the law \(\top\) can neither have an identity element nor be associative.
\item
  The right and left translations of a distributive quasi-group E are automorphisms of E.
\item
  If E is a finite distributive quasi-group, the structure induced on every stable subset of E is a distributive quasi-group structure.
\item
  Let E be a distributive quasi-group. If R is an equivalence relation which is left (resp. right) compatible (no. 3) with the law T, the classes mod. R are stable subsets of E. If E is finite, all these classes are obtained one from another by left (resp. right) translation; under the same conditions, if R is compatible with the law T, the quotient structure on E/R is a distributive quasi-group structure.
\item
  Let E be a distributive quasi-group. The set \(A_{a}\) of elements of E which are permutable with a given element a is stable; if E is finite, for all \(x \in A_{a}\) , \(A_{x} = A_{a}\) (note, using (e), that there exists \(y \in A_{a}\) such that \(x = a \top y\) ) and when x runs through E the sets \(A_{x}\) are identical with the equivalence classes with respect to a relation compatible with the law T.
\item
  If E is a finite distributive quasi-group and the law \(\top\) is commutative, the number of elements of E is odd (consider the ordered pairs \((x,y)\) of elements of E such that \(x\top y=y\top x=a\) , where a is given).
\end{enumerate}

\begin{enumerate}
\def\labelenumi{\arabic{enumi}.}
\setcounter{enumi}{6}
\tightlist
\item
  Suppose given on a set E an (associative and commutative) addition under which all the elements of E are invertible and an (associative) multiplication which is doubly distributive with respect to addition \(*\) (in other words, a ring structure) \(_{*}\) ; write \([x,y]=xy-yx\) ; the law \((x,y)\mapsto[x,y]\) is doubly distributive with respect to addition. For x and y to be permutable under multiplication, it is necessary and sufficient that \([x,y]=0\) ; and we have the identities
\end{enumerate}

\[
[ x, y ] = - [ y, x ]; \quad [ x, [ y, z ] ] + [ y, [ z, x ] ] + [ z, [ x, y ] ] = 0
\]

(the second is known as the ``Jacobi identity''). The second identity may also be written

\[
[ x, [ y, z ] ] - [ [ x, y ], z ] = [ [ z, x ], y ]
\]

which expresses the ``deviation from associativity'' of the law \([x, y]\) .

\begin{enumerate}
\def\labelenumi{\arabic{enumi}.}
\setcounter{enumi}{7}
\tightlist
\item
  With the same hypotheses as in Exercise 7, let \(x \top y = xy + yx\) ; then the law \(\top\) is commutative and doubly distributive with respect to addition but not in general associative.
\end{enumerate}

\begin{enumerate}
\def\labelenumi{(\alph{enumi})}
\item
  For all \(x \in \mathbb{E}\) , show that \(\top^{m+n} x = \left( \begin{array}{c} m \\ \top \end{array} x \right) \top \left( \begin{array}{c} n \\ \top \end{array} x \right)\) .
\item
  If we write \([x, y, z] = (x \top y) \top z - x \top (y \top z)\) (deviation from associativity of the law \(\top\) ), prove the identities
\end{enumerate}

\[
[ x, y, z ] + [ z, y, x ] = 0
\]

\[
[ x, y, z ] + [ y, z, x ] + [ z, x, y ] = 0
\]

\[
[ x \top y, u, z ] = [ u, [ (x \top y), z) ] ]
\]

(the notation \([x, y]\) meaning the same as in Exercise 7)

\[
[ x \top y, u, z ] + [ y \top z, u, x ] + [ z \top x, u, y ] = 0.
\]

¶ 9. Suppose given on E an (associative and commutative) addition under which all the elements of E are invertible and a multiplication which is non-associative, but commutative and doubly distributive with respect to addition. Suppose further that \(n \in \mathbf{Z}\) , \(n \neq 0\) and \(nx = 0\) imply \(x = 0\) in E. Show that if, writing \([x, y, z] = (xy)z - x(yz)\) , the identity

\[
[ x y, u, z ] + [ y z, u, z ] + [ z x, u, y ] = 0
\]

holds, then \(x^{m + n} = x^m x^n\) for all \(x\) (show, by induction on \(p\) , that the identity \([x^q, y, x^{p - q}] = 0\) holds for \(1 \leqslant q < p\) ).

\begin{enumerate}
\def\labelenumi{\arabic{enumi}.}
\setcounter{enumi}{9}
\tightlist
\item
  Let E be a commutative monoid whose law is written additively and whose identity element is denoted by 0. Let \(\top\) be an internal law on E which is distributive with respect to addition and such that \(0 \top x = x \top 0 = 0\) for all \(x \in \mathbb{E}\) . Let S be a subset of E which is stable under addition and under each of the external laws derived from \(\top\) ; let \(\overline{\mathbb{E}}\) denote the monoid of differences of E with respect to S under addition and \(\varepsilon\) the canonical homomorphism of E into \(\overline{\mathbb{E}}\) . Then there exists on \(\overline{\mathbb{E}}\) one and only one law \(\overline{\top}\) which is distributive with respect to addition and such that \(\varepsilon\) is a homomorphism for the laws \(\top\) and \(\overline{\top}\) . If the law \(\top\) is associative (resp. commutative), so is the law \(\overline{\top}\) .
\end{enumerate}

